\documentclass[10pt,a4paper]{article}

\usepackage[a4paper,margin=15mm,headheight=14pt]{geometry}
\usepackage{array}
\usepackage{booktabs}
\usepackage{caption}
\usepackage{enumitem}
\usepackage{fancyhdr}
\usepackage{float}
\usepackage{fontspec}
\usepackage{graphicx}
\usepackage{longtable}
\usepackage{microtype}
\usepackage{parskip}
\usepackage{tabularx}
\usepackage{xcolor}
\usepackage{hyperref}

\definecolor{ink}{HTML}{1F2937}
\definecolor{muted}{HTML}{64748B}
\definecolor{accent}{HTML}{006B3C}
\definecolor{secondary}{HTML}{0B7A46}
\definecolor{pale}{HTML}{EAF6EF}
\definecolor{panel}{HTML}{F5F7F6}
\definecolor{pending}{HTML}{FEF3C7}
\definecolor{line}{HTML}{CBD5E1}

\hypersetup{
  colorlinks=true,
  linkcolor=accent,
  urlcolor=secondary,
  pdftitle={TokTickIT Lab 2 Report Draft},
  pdfauthor={GUNTEE DOUNGMANEE},
}

\setmainfont{Latin Modern Roman}
\setlength{\parindent}{0pt}
\setlist{nosep,leftmargin=*}
\setlist[itemize]{label=\textcolor{accent}{\small\textbullet}}
\newcolumntype{Y}{>{\raggedright\arraybackslash}X}

\pagestyle{fancy}
\fancyhf{}
\lhead{TokTickIT Lab 2}
\rhead{Submission evidence draft}
\cfoot{\thepage}
\renewcommand{\headrulewidth}{0.3pt}

\newcommand{\captureplaceholder}[1]{%
  \fcolorbox{muted}{pending}{%
    \parbox[c][0.18\textheight][c]{0.93\linewidth}{%
      \centering\textcolor{ink}{\textbf{Evidence pending}}\\[3pt]
      \small #1
    }%
  }%
}

\newcommand{\identitycard}[5]{%
  \fcolorbox{accent}{panel}{%
    \parbox[t]{\dimexpr\linewidth-2\fboxsep-2\fboxrule\relax}{%
      {\small\bfseries\textcolor{accent}{#1}}\par
      {\bfseries #2}\par
      {\small ID: #3}\par
      {\small GitHub: \href{#4}{#5}}%
    }%
  }%
}

\newcommand{\evidenceimage}[2]{%
  \IfFileExists{#1}{%
    \includegraphics[width=\linewidth,height=0.18\textheight,keepaspectratio]{#1}%
  }{\captureplaceholder{#2}}%
}

% Long browser-use pages are kept intact in evidence but cropped in the PDF
% presentation so that the relevant labels remain readable.
\newcommand{\evidenceimagecrop}[3]{%
  \IfFileExists{#1}{%
    \includegraphics[width=\linewidth,height=0.26\textheight,keepaspectratio,trim=#2,clip]{#1}%
  }{\captureplaceholder{#3}}%
}

\newcommand{\workflowimage}[2]{%
  \IfFileExists{#1}{%
    \includegraphics[width=\linewidth,height=0.25\textheight,keepaspectratio]{#1}%
  }{\captureplaceholder{#2}}%
}

\newcommand{\twocaptures}[6]{%
\begin{figure}[H]
  \centering
  \begin{minipage}[t]{0.48\linewidth}
    \centering
    \evidenceimage{#1}{#2}
    \captionof{figure}{#3}
  \end{minipage}\hfill
  \begin{minipage}[t]{0.48\linewidth}
    \centering
    \evidenceimage{#4}{#5}
    \captionof{figure}{#6}
  \end{minipage}
\end{figure}
}

\newcommand{\twocapturescrop}[8]{%
\begin{figure}[H]
  \centering
  \begin{minipage}[t]{0.48\linewidth}
    \centering
    \evidenceimagecrop{#1}{#2}{#3}
    \captionof{figure}{#4}
  \end{minipage}\hfill
  \begin{minipage}[t]{0.48\linewidth}
    \centering
    \evidenceimagecrop{#5}{#6}{#7}
    \captionof{figure}{#8}
  \end{minipage}
\end{figure}
}

\newcommand{\widecapturecrop}[4]{%
\begin{figure}[H]
  \centering
  \IfFileExists{#1}{%
    \includegraphics[width=0.92\linewidth,height=0.30\textheight,keepaspectratio,trim=#2,clip]{#1}%
  }{\captureplaceholder{#3}}%
  \caption{#4}
\end{figure}
}

\newcommand{\threecapture}[9]{%
\begin{figure}[H]
  \centering
  \begin{minipage}[t]{0.31\linewidth}
    \centering
    \evidenceimage{#1}{#2}
    \captionof{figure}{#3}
  \end{minipage}\hfill
  \begin{minipage}[t]{0.31\linewidth}
    \centering
    \evidenceimage{#4}{#5}
    \captionof{figure}{#6}
  \end{minipage}\hfill
  \begin{minipage}[t]{0.31\linewidth}
    \centering
    \evidenceimage{#7}{#8}
    \captionof{figure}{#9}
  \end{minipage}
\end{figure}
}

\newcommand{\statusnote}[1]{%
  \fcolorbox{accent}{pale}{\parbox{0.96\linewidth}{\small #1}}%
}

\begin{document}

\begin{titlepage}
  \thispagestyle{empty}
  \vspace*{15mm}
  {\Large\textcolor{accent}{CPE 334 \textbar\ Lab 2}}\par
  \vspace{8mm}
  {\Huge\bfseries TokTickIT}\par
  \vspace{3mm}
  {\Large Requester-facing ticket workflow}\par
  \vfill
  \begin{tabular}{@{}>{\bfseries}p{42mm}p{105mm}@{}}
    Student & GUNTEE DOUNGMANEE \\
    Student ID & 67070501003 \\
    Peer reviewer & POLWARIT WATTHANAHEMMARAT (MadMax168), ID 67070501067 \\
    Repository & \href{https://github.com/ovenmakemeheat/toktickit}{ovenmakemeheat/toktickit} \\
    Report Issue & \href{https://github.com/ovenmakemeheat/toktickit/issues/65}{Issue \#65 - Lab 2 final submission report} \\
    Integrated baseline & \texttt{lab2-staging} at \texttt{f6d00cb} \\
    Stable branch & \texttt{main} at \texttt{538b5da} (Lab 2 release) \\
    Stack & React, TypeScript, Vite, Bootstrap, Express, Prisma, PostgreSQL \\
  \end{tabular}
  \vfill
  \statusnote{\textbf{DRAFT - 2026-09-05.} Lab 2 implementation is integrated on \texttt{lab2-staging} and released to final \texttt{main}. Eight UI report-state captures, seven final-main \texttt{browser-use} captures, final-main client/server test output, the final-main-aligned E2E rerun, and the current Project board capture are recorded. Human visual sign-off and formal GitHub \texttt{APPROVED} review evidence remain pending.}
  \vfill
  {\small One concise PDF is required by the labsheet; this source is the LaTeX draft for that PDF.}
\end{titlepage}

\tableofcontents
\clearpage

\section{Answer Part 1: Git Use with Engineering Workflow}

Lab 2 uses one feature branch per Issue and integrates each reviewed Pull
Request into \texttt{lab2-staging}. The human reviewer merges the feature PR;
the PR author does not merge it. The current repository state is:

\begin{verbatim}
main (538b5da, Lab 2 release; PR #70)
|-- lab2-staging (f6d00cb, integrated Lab 2 baseline)
    |-- feature/5-sprint-specification   (#51 -> PR #58)
    |-- feature/6-requester-context     (#52 -> PR #59)
    |-- feature/7-create-ticket         (#53 -> PR #60)
    |-- feature/8-my-tickets            (#54 -> PR #61)
    |-- feature/9-ticket-detail         (#55 -> PR #62)
    |-- feature/10-zen-green            (#56 -> PR #63)
    |-- feature/11-e2e-evidence          (#57 -> PR #64)
    |-- feature/12-report-compliance     (#65 -> PR #66)
    `-- feature/13-test-isolation        (#68 -> PR #69)
`-- docs/lab2-docs (current evidence branch, based on main)
\end{verbatim}

\begin{center}
\small
\begin{tabular}{@{}p{12mm}p{49mm}p{24mm}p{66mm}@{}}
  \toprule
  Issue & Feature branch & PR & Current state \\
  \midrule
  \href{https://github.com/ovenmakemeheat/toktickit/issues/51}{\#51} & \path{feature/5-sprint-specification} & \href{https://github.com/ovenmakemeheat/toktickit/pull/58}{\#58} & Merged into \texttt{lab2-staging} by \texttt{MadMax168}. \\
  \href{https://github.com/ovenmakemeheat/toktickit/issues/52}{\#52} & \path{feature/6-requester-context} & \href{https://github.com/ovenmakemeheat/toktickit/pull/59}{\#59} & Merged into \texttt{lab2-staging} by \texttt{MadMax168}. \\
  \href{https://github.com/ovenmakemeheat/toktickit/issues/53}{\#53} & \path{feature/7-create-ticket} & \href{https://github.com/ovenmakemeheat/toktickit/pull/60}{\#60} & Merged into \texttt{lab2-staging} by \texttt{MadMax168}. \\
  \href{https://github.com/ovenmakemeheat/toktickit/issues/54}{\#54} & \path{feature/8-my-tickets} & \href{https://github.com/ovenmakemeheat/toktickit/pull/61}{\#61} & Merged into \texttt{lab2-staging} by \texttt{MadMax168}. \\
  \href{https://github.com/ovenmakemeheat/toktickit/issues/55}{\#55} & \path{feature/9-ticket-detail} & \href{https://github.com/ovenmakemeheat/toktickit/pull/62}{\#62} & Merged into \texttt{lab2-staging} by \texttt{MadMax168}. \\
  \href{https://github.com/ovenmakemeheat/toktickit/issues/56}{\#56} & \path{feature/10-zen-green} & \href{https://github.com/ovenmakemeheat/toktickit/pull/63}{\#63} & Merged into \texttt{lab2-staging} by \texttt{MadMax168}. \\
  \href{https://github.com/ovenmakemeheat/toktickit/issues/57}{\#57} & \path{feature/11-e2e-evidence} & \href{https://github.com/ovenmakemeheat/toktickit/pull/64}{\#64} & Merged into \texttt{lab2-staging} by \texttt{MadMax168}. \\
  \href{https://github.com/ovenmakemeheat/toktickit/issues/65}{\#65} & \path{feature/12-report-compliance} & \href{https://github.com/ovenmakemeheat/toktickit/pull/66}{\#66} & Merged into \texttt{lab2-staging} by \texttt{MadMax168}. \\
  \href{https://github.com/ovenmakemeheat/toktickit/issues/68}{\#68} & \path{feature/13-test-isolation} & \href{https://github.com/ovenmakemeheat/toktickit/pull/69}{\#69} & Merged into \texttt{lab2-staging} by \texttt{MadMax168}. \\
  Release & \path{lab2-staging} & \href{https://github.com/ovenmakemeheat/toktickit/pull/70}{\#70} & Merged into final \texttt{main} by \texttt{MadMax168}. \\
  \bottomrule
\end{tabular}
\end{center}

\textbf{Student and peer reviewer}

\vspace{1mm}
\begin{tabular}{@{}p{0.47\linewidth}@{\hspace{4mm}}p{0.47\linewidth}@{}}
  \identitycard{Student}{GUNTEE DOUNGMANEE}{67070501003}{https://github.com/ovenmakemeheat}{ovenmakemeheat}
  &
  \identitycard{Peer reviewer}{POLWARIT WATTHANAHEMMARAT}{67070501067}{https://github.com/MadMax168}{MadMax168}
\end{tabular}

\vspace{3mm}
\textbf{Rendered \texttt{reviewer.md}}

The reviewer record identifies the same student and peer reviewer used for Lab 1.
It records the Issue-to-PR links, reviewer comments, author responses, merge
accounts, and the two delivery items that still require human action: formal
GitHub \texttt{APPROVED} events and final visual sign-off.

{\scriptsize
\begin{longtable}{@{}p{13mm}p{52mm}p{65mm}p{25mm}@{}}
  \toprule
  PR & Human review comment or decision & Author response / resulting change & State \\
  \midrule
  \endfirsthead
  \toprule
  PR & Human review comment or decision & Author response / resulting change & State \\
  \midrule
  \endhead
  \href{https://github.com/ovenmakemeheat/toktickit/pull/58}{\#58} & \texttt{MadMax168}: ``ready to merge, Approve it.'' & No change request was recorded before merge. & Merged; formal approval pending \\
  \href{https://github.com/ovenmakemeheat/toktickit/pull/59}{\#59} & Requested a UI loading-state test and responsive selector coverage; then said ``All well, Ready to merge.'' & Added requester-loading coverage and responsive selector assertions in commit \texttt{68838e2}. & Addressed and merged \\
  \href{https://github.com/ovenmakemeheat/toktickit/pull/60}{\#60} & \texttt{MadMax168} reported no bug and approved the merge in a PR comment. & No change request was recorded. & Merged; formal approval pending \\
  \href{https://github.com/ovenmakemeheat/toktickit/pull/61}{\#61} & Requested protection against stale/out-of-order list responses and restoration of the visible Open Ticket action. & Added \texttt{AbortController}, a latest-request guard, and a deferred-response regression test in commit \texttt{6719a96}; restored the visible action. & Addressed and merged \\
  \href{https://github.com/ovenmakemeheat/toktickit/pull/62}{\#62} & Requested ownership/context validation before multipart parsing and enforcement of the active-attachment limit. & Addressed both findings in commit \texttt{c0a410a}, including the cross-owner multipart regression. & Addressed and merged \\
  \href{https://github.com/ovenmakemeheat/toktickit/pull/63}{\#63} & \texttt{MadMax168}: ``Good, Ready to merge.'' & No change request was recorded. & Merged; formal approval pending \\
  \href{https://github.com/ovenmakemeheat/toktickit/pull/64}{\#64} & Requested deterministic Ticket Number collision-retry coverage, corrected test traceability, and complete responsive evidence. & Added collision-retry coverage and corrected the evidence/count records in commit \texttt{9380c9c}; the author replied to the discussion. & Addressed and merged \\
  \bottomrule
\end{longtable}
}

\emph{The feature PRs were merged by \texttt{MadMax168}, not by the PR author.
The record does not convert ready-to-merge comments into formal approvals.}

The complete review index and response summaries are in
\href{https://github.com/ovenmakemeheat/toktickit/blob/main/docs/lab-02/reviewer.md}{\texttt{docs/lab-02/reviewer.md}}. The record identifies
\texttt{MadMax168} as the human reviewer and records the requested changes and
author responses for PRs \#59, \#61, \#62, and \#64. GitHub currently exposes
human ready-to-merge or approve comments, but its formal \texttt{reviewDecision}
is empty for PRs \#58-\#64; the report must not label those comments as formal
\texttt{APPROVED} reviews. The final release PR \#70 records the merge from
\texttt{lab2-staging} into \texttt{main}. The Lab 2 Kanban view was captured
with the \texttt{browser-use} CLI at
\texttt{/users/ovenmakemeheat/projects/2/views/2}; it shows the six Lab 1-aligned
columns (\texttt{Backlog}, \texttt{Specified}, \texttt{Started}, \texttt{PR
Review}, \texttt{Fixing}, and \texttt{Done}) with Issues \#51-\#57 in
\texttt{Done}.

\textbf{Required Part 1 captures}

\widecapturecrop{docs/lab-02/report/evidence/final-main-history.png}{210pt 1450pt 210pt 0pt}{Capture final-main history.}{Final \texttt{main} history showing the latest release, staging, and Lab 2 commits (cropped for legibility).}
\widecapturecrop{docs/lab-02/report/evidence/final-lab2-project-board.png}{0pt 0pt 0pt 0pt}{Capture the final Lab 2 Kanban board.}{Final Lab 2 Kanban board with the six Lab 1-aligned workflow columns; Issues \#51-\#57 are in \texttt{Done}.}

\twocapturescrop{docs/lab-02/report/evidence/final-main-readme.png}{210pt 650pt 210pt 0pt}{Render final-main README.}{Repository README setup evidence (cropped to the readable content area).}{docs/lab-02/report/evidence/final-main-gitignore.png}{210pt 0pt 210pt 0pt}{Render final-main .gitignore.}{Secrets, dependencies, build output, and LaTeX artifact exclusions.}

\widecapturecrop{docs/lab-02/report/evidence/final-main-repository-tree.png}{210pt 1770pt 210pt 0pt}{Capture the final main repository tree.}{Final \texttt{main} repository tree captured with \texttt{browser-use} (cropped to the tree and README entry).}

\textbf{Repository directory structure}

The final-main repository tree follows the directory-structure evidence format
used in the Lab 1 report. The Lab 2 paths visible in the repository are:

\begin{verbatim}
docs/lab-02/        specification.md  tests.md  ui-spec.md  api-spec.md
                   reviewer.md  ai-use.md  report/
server/tests/lab-02/  create-ticket.api.test.ts  my-tickets.api.test.ts
                      ticket-detail.api.test.ts  attachments.api.test.ts
client/tests/lab-02/  CreateTicket.test.tsx  MyTickets.test.tsx
                      RequesterTicketDetail.test.tsx  AttachmentSection.test.tsx
e2e/lab-02/           requester-ticket-flow.spec.ts  report-evidence.spec.ts
artifacts/lab-02/screenshots/  create-ticket/  my-tickets/  ticket-detail/
\end{verbatim}

The final-main tree, README, \texttt{.gitignore}, history, and Project board
captures were taken from GitHub with \texttt{browser-use}. The Kanban board
capture records the six Lab 1-aligned columns and the seven Lab 2 scope Issues
as \texttt{Done}; human visual sign-off remains a separate review
responsibility. The original full
browser captures remain in the evidence directory; the long pages are
content-focused crops in this PDF for legibility.

\section{Answer Part 2: Spec DD}

The complete contract is linked at
\href{https://github.com/ovenmakemeheat/toktickit/blob/main/docs/lab-02/specification.md}{\texttt{docs/lab-02/specification.md}} and is cross-referenced by the
API, UI, and test contracts. It contains 19 functional requirements, 26
business rules, 24 acceptance criteria, and a Definition of Done. The
contract-first document was delivered in PR \#58 before the implementation PRs
were completed.

GitHub records the contract commit as
\href{https://github.com/ovenmakemeheat/toktickit/commit/8f372bd5f7dc687686b61273817c30e9ce24d14b}{\texttt{8f372bd}}
(2026-08-24 06:30 UTC). PR \#58 was opened at 06:31 UTC and merged into
\texttt{lab2-staging} at 09:14 UTC; the next implementation PR, PR \#59, was
opened later that day. The browser-use capture below shows PR \#58's merged
state, source branch, contract-only summary, commit \texttt{8f372bd}, and merge
into \texttt{lab2-staging}; this supports the contract-before-implementation
sequence. The final release PR \#70 was subsequently merged into
\texttt{main} at \texttt{538b5da}. The final-main browser-use captures below
show the released contract and merge timeline.

\begin{center}
\small
\begin{tabular}{@{}p{42mm}p{55mm}p{59mm}@{}}
  \toprule
  Contract area & Selected decision & Traceability source \\
  \midrule
  Requester context & Development Requester selection is a test context, not authentication. & FR-01--FR-04, BR-03--BR-07, AC-01--AC-03 \\
  Ticket creation & The server owns Ticket Number, Ticket Date, requester ownership, and initial \texttt{NEW} status. & FR-06--FR-08, BR-01--BR-13, AC-05--AC-10 \\
  Ownership & List, detail, and attachment operations are scoped to the selected Requester. & FR-09--FR-16, BR-23, AC-11--AC-21 \\
  Scope boundary & Authentication, IT Staff workflow, collaboration, and later status transitions remain excluded. & Scope and DoD sections \\
  \bottomrule
\end{tabular}
\end{center}

\textbf{Rendered \texttt{specification.md} index}

\begin{center}
\small
\begin{tabular}{@{}p{42mm}p{112mm}@{}}
  \toprule
  Contract section & Numbered content rendered in the final-main contract \\
  \midrule
  Functional requirements & FR-01 through FR-19: requester context, ticket creation, owned list/detail, attachments, errors, and responsive accessibility. \\
  Business rules & BR-01 through BR-26: server-owned values, validation, ownership, query behavior, attachment policy, and explicit state handling. \\
  Acceptance criteria & AC-01 through AC-24: selection, creation, list/detail, attachment lifecycle, ownership, responsive behavior, and safe errors. \\
  Definition of Done & Product completion, traceable passing tests, contract conformance, final-main verification, peer review, visual review, and submission material. \\
  \bottomrule
\end{tabular}
\end{center}

\textbf{Rendered contract evidence}

\statusnote{The source contract is complete, implemented, and present on final \texttt{main}. The browser-use captures below provide the rendered final-main contract and the merged release PR timeline.}

\twocapturescrop{docs/lab-02/report/evidence/final-main-release-pr.png}{210pt 0pt 210pt 0pt}{Capture final release PR.}{PR \#70 merged \texttt{lab2-staging} into final \texttt{main}, including the test-isolation release.}{docs/lab-02/report/evidence/final-main-specification.png}{210pt 5660pt 210pt 0pt}{Render final-main specification overview.}{Rendered \texttt{specification.md} overview with numbered FR, BR, AC, and DoD content (cropped for legibility).}

\section{Answer Part 3: Test DD and Traceability}

The planned matrix and acceptance-criteria mapping are maintained in
\href{https://github.com/ovenmakemeheat/toktickit/blob/main/docs/lab-02/tests.md}{\texttt{docs/lab-02/tests.md}}. The current final-main validation
record contains the following results:

\textbf{Rendered \texttt{tests.md}: planned tests and actual paths}

{\scriptsize
\begin{longtable}{@{}p{21mm}p{24mm}p{89mm}p{21mm}@{}}
  \toprule
  ID & Boundary & Actual test-file path or evidence path & Result \\
  \midrule
  \endfirsthead
  \toprule
  ID & Boundary & Actual test-file path or evidence path & Result \\
  \midrule
  \endhead
  UNIT-01 & Unit & \path{server/tests/lab-02/ticket-number.unit.test.ts} & Passed \\
  UNIT-02 & Unit & \path{server/tests/lab-02/ticket-validation.unit.test.ts} & Passed \\
  UNIT-03 & Unit & \path{server/tests/lab-02/ticket-query.unit.test.ts} & Passed \\
  UNIT-04 & Unit & \path{server/tests/lab-02/attachment-policy.unit.test.ts} & Passed \\
  \mbox{API-01--API-04} & API & \path{server/tests/lab-02/create-ticket.api.test.ts} & Passed \\
  API-05 & API & \path{server/tests/lab-02/attachments.api.test.ts} & Passed \\
  \mbox{API-06--API-08} & API & \path{server/tests/lab-02/my-tickets.api.test.ts} & Passed \\
  \mbox{API-09--API-10} & API & \path{server/tests/lab-02/ticket-detail.api.test.ts} & Passed \\
  \mbox{API-11--API-15} & API & \path{server/tests/lab-02/attachments.api.test.ts} & Passed \\
  API-16 & API & \path{server/tests/lab-02/create-ticket.api.test.ts}; \path{server/tests/lab-02/attachments.api.test.ts} & Passed \\
  UI-01--UI-02 & UI & \path{client/tests/lab-02/RequesterSelection.test.tsx} & Passed \\
  UI-03--UI-05 & UI & \path{client/tests/lab-02/CreateTicket.test.tsx} & Passed \\
  UI-06--UI-07 & UI & \path{client/tests/lab-02/MyTickets.test.tsx} & Passed \\
  UI-08 & UI & \path{client/tests/lab-02/RequesterTicketDetail.test.tsx} & Passed \\
  UI-09--UI-11 & UI & \path{client/tests/lab-02/AttachmentSection.test.tsx} & Passed \\
  UI-12 & UI & \path{client/tests/lab-02/ZenGreenResponsive.test.tsx} & Passed \\
  STYLE-01 & Style & \path{client/tests/lab-02/ZenGreenResponsive.test.tsx} & Passed \\
  RESP-01 & Responsive & \path{e2e/lab-02/requester-ticket-flow.spec.ts} & Passed \\
  VIS-01 & Visual & \path{artifacts/lab-02/screenshots/}; \path{docs/lab-02/report/evidence/} & Evidence recorded; sign-off pending \\
  E2E-01--E2E-02 & E2E & \path{e2e/lab-02/requester-ticket-flow.spec.ts} & Passed \\
  \bottomrule
\end{longtable}
}

\textbf{Acceptance-criterion traceability rendered from \texttt{tests.md}}

{\scriptsize
\begin{tabular}{@{}p{13mm}p{66mm}p{13mm}p{66mm}@{}}
  \toprule
  AC & Planned test IDs & AC & Planned test IDs \\
  \midrule
  AC-01 & UI-01, E2E-01 & AC-13 & API-08, UI-06 \\
  AC-02 & API-01, UI-01, E2E-01 & AC-14 & API-08, UI-07 \\
  AC-03 & UI-02, E2E-01 & AC-15 & API-09, UI-08, E2E-01 \\
  AC-04 & API-01, UI-03, E2E-01 & AC-16 & API-10, UI-08 \\
  AC-05 & UNIT-01, API-02, UI-03, E2E-01 & AC-17 & API-11, UI-09, E2E-02 \\
  AC-06 & UNIT-02, UI-04 & AC-18 & UNIT-04, API-05, API-12, UI-10, E2E-02 \\
  AC-07 & UNIT-02, API-03 & AC-19 & API-13, UI-11, E2E-02 \\
  AC-08 & API-04, UI-05 & AC-20 & UNIT-04, API-14, UI-11, E2E-02 \\
  AC-09 & API-05, E2E-02 & AC-21 & API-10, API-13, API-15, UI-11, E2E-02 \\
  AC-10 & API-05, UI-05, E2E-02 & AC-22 & STYLE-01, RESP-01, VIS-01, E2E-01, E2E-02 \\
  AC-11 & API-06, UI-06, E2E-01 & AC-23 & UI-12, STYLE-01, VIS-01 \\
  AC-12 & UNIT-03, API-07, UI-06, E2E-01 & AC-24 & API-16 \\
  \bottomrule
\end{tabular}
}

\begin{center}
\small
\begin{tabular}{@{}p{47mm}p{28mm}p{74mm}@{}}
  \toprule
  Boundary & Result & Main evidence path \\
  \midrule
  Client unit/UI suite & 8 files, 32 tests passed & \texttt{client/tests/lab-01/} and \texttt{client/tests/lab-02/} \\
  Server unit/API suite & 13 files, 70 tests passed & \texttt{server/tests/lab-01/} and \texttt{server/tests/lab-02/} \\
  Combined test command & 21 files, 102 tests passed & \texttt{bun run test} \\
  Browser E2E & 2 tests passed and 44 screenshots & \path{e2e/lab-02/} \\
  Full repository gate & Passed on final-main baseline & \texttt{bun run verify}; Biome, hooks, Prisma, types, tests, and builds \\
  \bottomrule
\end{tabular}
\end{center}

The principal Lab 2 test paths are:

\begin{itemize}
  \item \path{server/tests/lab-02/create-ticket.api.test.ts}, \path{server/tests/lab-02/my-tickets.api.test.ts}, \path{server/tests/lab-02/ticket-detail.api.test.ts}, and \path{server/tests/lab-02/attachments.api.test.ts};
  \item \path{client/tests/lab-02/CreateTicket.test.tsx}, \path{client/tests/lab-02/MyTickets.test.tsx}, \path{client/tests/lab-02/RequesterTicketDetail.test.tsx}, and \path{client/tests/lab-02/AttachmentSection.test.tsx}.
\end{itemize}

Every AC is mapped to at least one test ID in the Markdown matrix.

The original lifecycle test produced 36 responsive screenshots. The dedicated
report-evidence test at \path{e2e/lab-02/report-evidence.spec.ts} produced eight
additional desktop captures for requester selection, Create Ticket states, and
the combined My Tickets filter/sort/page state.

The 2026-09-04 elevated local run of \texttt{bun run verify} passed on the
\texttt{docs/lab2-docs} branch, which points to final \texttt{main} at
\texttt{538b5da}. It checked
Biome, Lefthook, Prisma, client and server type checking, 21 test files with
102 tests, and both workspace builds. The normal sandbox invocation stopped at
Prisma dependency access with an OS \texttt{EPERM}; this is recorded as an
environment limitation, not as a passing test result. The supplied final-main
test capture separately shows the client and server test commands completing
with 8 files/32 tests and 13 files/70 tests.

\begin{figure}[H]
  \centering
  \includegraphics[height=0.23\textheight,keepaspectratio]{docs/lab-02/report/evidence/final-main-test-output.png}
  \caption{Final-main client and server test output captured after the Lab 2 release (resized for the report).}
\end{figure}

\section{Answer Part 4: AI Use with Reflection}

The AI-use record is linked at
\href{https://github.com/ovenmakemeheat/toktickit/blob/main/docs/lab-02/ai-use.md}{\texttt{docs/lab-02/ai-use.md}}. The implementation assistant was OpenAI
Codex using GPT-5.6 Luna max. The student reviewed the PDF-derived contracts, generated
changes, tests, screenshots, and review responses.

\begin{center}
\small
\begin{longtable}{@{}p{9mm}p{65mm}p{77mm}@{}}
  \toprule
  ID & Selected prompt focus & Effect \\
  \midrule
  P1 & Read the Lab 2 PDF and convert workflows and submission parts into a contract. & Established numbered FR, BR, AC, and DoD traceability. \\
  P2 & Create the Issue and branch plan from the contract. & Kept one feature branch per Lab 2 Issue and preserved the staging flow. \\
  P3 & Implement data, seed, requester context, API, and ownership boundaries. & Added the requester-facing backend contract without authentication or later-lab scope. \\
  P4 & Implement Create Ticket, My Tickets, Detail, and attachments. & Connected the four product areas to the approved API and UI decisions. \\
  P5 & Add tests with explicit FR/BR/AC traceability. & Produced unit, API, client UI, responsive, and E2E coverage. \\
  P6 & Inspect the requester flow at desktop, tablet, and mobile sizes. & Generated the responsive screenshot inventory and overflow assertions. \\
  P7 & Review the complete Lab 2 PR sequence and address comments. & Updated implementation, tests, evidence, and reviewer records after human feedback. \\
  P8 & Check report readiness against Answer Parts 1-9. & Separated staging evidence from final-main and missing-capture work. \\
  \bottomrule
\end{longtable}
\end{center}

\textbf{My Reflection}

AI assistance was most useful for translating the labsheet into small,
traceable contracts and checking the relationships between implementation,
tests, evidence, and review records. Human responsibility remained essential:
the student confirmed scope, reviewed generated changes, inspected screenshots,
validated the environment, and did not treat automated review output as peer
approval.

\section{Answer Part 5: Development Requester Select Screen}

Lab 2 uses a simulated login screen to select an active Development Requester
for the session. It is explicitly not authentication. The repeat-safe seed
contains active Requesters A-D and one inactive Requester; the selector returns
only the four active records. Continue is disabled until a valid active
selection is made, and Change Requester clears the previous requester-specific
view.

\twocaptures{docs/lab-02/report/evidence/requester-selection-ready.png}{Requester selection before a choice.}{Active requester options with Continue disabled and inactive Requesters excluded.}{docs/lab-02/report/evidence/requester-selection-selected.png}{Requester selection after a choice.}{Requester A selected and Continue enabled.}

\twocaptures{docs/lab-02/report/evidence/requester-summary.png}{Selected requester context.}{Selected Requester A shown in the requester summary.}{docs/lab-02/report/evidence/requester-change.png}{Change Requester behavior.}{Change Requester returns to a cleared selection screen.}

The four captures were collected by the dedicated Playwright evidence test
using the seeded active requester shape. They show the selection, enabled
continuation, selected-user summary, and reset behavior without presenting this
simulated context as authentication.

\section{Answer Part 6: Working Ticket Screen - Create Mode}

Create Ticket accepts Category, Related System, Requested Priority, Summary,
Description, and permitted attachments. The selected requester is displayed as
read-only context. The server generates the official Ticket Number, Ticket Date,
requester ownership, and initial \texttt{NEW} status. Attachment upload is a
separate operation so a saved Ticket remains visible if an upload fails.

\twocaptures{docs/lab-02/report/evidence/requester-summary.png}{Selected requester context before Create Ticket.}{Requester selection and matching context.}{artifacts/lab-02/screenshots/create-ticket/create-ticket-reference-loading-desktop.png}{Capture the reference-data loading state.}{Reference data loading from the API.}

\twocaptures{docs/lab-02/report/evidence/create-ticket-initial-desktop.png}{Dedicated initial-form capture.}{Initial editable form with required indicators.}{artifacts/lab-02/screenshots/create-ticket/create-ticket-validation-desktop.png}{Existing validation evidence.}{Field-level messages for an invalid submission.}

\twocaptures{docs/lab-02/report/evidence/create-ticket-submitting-desktop.png}{Dedicated submitting-state capture.}{Disabled Submitting button and duplicate-submit protection.}{artifacts/lab-02/screenshots/create-ticket/create-ticket-success-desktop.png}{Existing success evidence.}{Backend-generated Ticket Number and saved values.}

\twocaptures{artifacts/lab-02/screenshots/create-ticket/create-ticket-api-failure-desktop.png}{Existing API failure evidence.}{Safe API failure with entered values preserved.}{docs/lab-02/report/evidence/create-ticket-invalid-attachment.png}{Dedicated mixed-attachment capture.}{One permitted PDF and one rejected file are shown before submission.}

The current UI/API tests and E2E evidence support the following contract
claims:

\begin{itemize}
  \item API-02 and E2E-01 verify a successful create response, saved Ticket, backend Ticket Number, and requester context.
  \item UI-04 and API-03 verify field-level validation and no invalid Ticket persistence.
  \item UI-05 and API-05 verify busy-submit protection, preserved values, and separate attachment failure behavior.
  \item The create API test verifies response \texttt{requester.id} and then checks the saved database \texttt{requesterId}, \texttt{categoryId}, \texttt{relatedSystemId}, and \texttt{currentStatus=NEW}; the response intentionally does not expose a raw \texttt{requesterId} field.
\end{itemize}

\section{Answer Part 7: Working My Tickets Screen}

My Tickets is always scoped by the selected Development Requester. It supports
search over Ticket Number and Summary, Category/Related System/Priority/Status
filters, deterministic sorting, one-based pagination, empty and no-results
states, and recoverable API failure.

\twocaptures{artifacts/lab-02/screenshots/my-tickets/my-tickets-filtered-desktop.png}{Existing filtered-ticket evidence.}{Requester-owned results with list controls.}{artifacts/lab-02/screenshots/my-tickets/my-tickets-ownership-isolation-desktop.png}{Existing ownership evidence.}{A different Requester cannot see the original Requester's Tickets.}

\twocaptures{artifacts/lab-02/screenshots/my-tickets/my-tickets-empty-desktop.png}{Existing empty-list evidence.}{No Tickets for the current Requester.}{artifacts/lab-02/screenshots/my-tickets/my-tickets-api-failure-desktop.png}{Existing API-failure evidence.}{Recoverable list failure.}

\begin{figure}[H]
  \centering
  \workflowimage{docs/lab-02/report/evidence/my-tickets-filter-sort-page-desktop.png}{Capture the combined filter, sort, and page transition state.}
  \caption{My Tickets with a non-default filter, ascending sort, and page 2 of 2.}
\end{figure}

The client and server tests verify search, filters, sorting, pagination,
ownership isolation, empty/no-results distinction, stale-response protection,
and API failure. The additional report capture makes the non-default Category,
Summary sort, ascending direction, and page 2 transition visible together.
It is a controlled UI-state fixture for the query controls; the final-main
browser captures document the released source and test baseline.

\section{Answer Part 8: Working Ticket Screen - View Mode and Attachments}

Owned Ticket Detail renders Ticket fields read-only and separates the
Attachment section. Active attachments can be uploaded and downloaded. Removal
requires confirmation and a 3-200 character reason, retains metadata, and
returns \texttt{410 ATTACHMENT\_REMOVED} when a removed file is requested.
Cross-Requester Ticket and Attachment access is rejected safely.

\twocaptures{artifacts/lab-02/screenshots/ticket-detail/ticket-detail-active-desktop.png}{Existing owned-detail evidence.}{Read-only Ticket Detail with active attachment.}{artifacts/lab-02/screenshots/ticket-detail/ticket-detail-invalid-attachment-desktop.png}{Existing invalid-attachment evidence.}{Unsupported attachment feedback.}

\twocaptures{artifacts/lab-02/screenshots/ticket-detail/ticket-detail-removed-desktop.png}{Existing removal evidence.}{Retained removed metadata and reason.}{artifacts/lab-02/screenshots/ticket-detail/ticket-detail-blocked-download-desktop.png}{Existing blocked-download evidence.}{Removed content is not downloadable or previewable.}

API-09 through API-16 and E2E-02 cover owned detail, attachment metadata,
upload, active download, soft removal, repeat removal, invalid files, storage
errors, removed-download rejection, and cross-owner access. The final report
should include a readable test-output or API excerpt for the unauthorized
Ticket/Attachment requests in addition to the UI screenshots. The relevant
test assertions are summarized here:

\begin{verbatim}
GET /api/tickets/<owned-ticket-id> [Requester B] -> 404 TICKET_NOT_FOUND
GET /api/tickets/<owned-ticket-id>/attachments [Requester B] -> 404 TICKET_NOT_FOUND
POST /api/tickets/<owned-ticket-id>/attachments [Requester B] -> 404 TICKET_NOT_FOUND
GET /api/tickets/<owned-ticket-id>/attachments/<id>/download [Requester B] -> 404 TICKET_NOT_FOUND
DELETE /api/tickets/<owned-ticket-id>/attachments/<id> [Requester B] -> 404 TICKET_NOT_FOUND
\end{verbatim}

These are asserted by \path{server/tests/lab-02/ticket-detail.api.test.ts} and
\path{server/tests/lab-02/attachments.api.test.ts}; no Ticket or Attachment
summary is disclosed to the other Requester.

\section{Answer Part 9: Zen Green UI and Responsive Evidence}

The complete visual contract is linked at
\href{https://github.com/ovenmakemeheat/toktickit/blob/main/docs/lab-02/ui-spec.md}{\texttt{docs/lab-02/ui-spec.md}}. The final-main-aligned E2E flow checks
1280x900, 820x900, and 390x844 viewports. The current inventory contains 36
responsive lifecycle captures plus eight desktop report-state captures. The
final-main evidence set adds seven \texttt{browser-use} captures and one
client/server test-output capture.

\widecapturecrop{docs/lab-02/report/evidence/final-main-ui-specification.png}{210pt 3970pt 210pt 0pt}{Render final-main UI specification overview.}{Rendered final-main \texttt{ui-spec.md} overview captured with \texttt{browser-use} (cropped for legibility).}

\threecapture{artifacts/lab-02/screenshots/create-ticket/create-ticket-validation-desktop.png}{Capture a readable desktop state.}{Desktop validation.}{artifacts/lab-02/screenshots/my-tickets/my-tickets-filtered-tablet.png}{Capture a readable tablet state.}{Tablet list controls.}{artifacts/lab-02/screenshots/ticket-detail/ticket-detail-blocked-download-mobile.png}{Capture a readable mobile state.}{Mobile blocked-download state.}

\begin{center}
\small
\begin{tabular}{@{}p{59mm}p{35mm}p{62mm}@{}}
  \toprule
  Visual checklist item & Integrated evidence & Final report status \\
  \midrule
  Zen Green colors and shared components & Client contract tests and screenshots & Rendered final-main \texttt{ui-spec.md}; human sign-off pending \\
  Editable/read-only fields & Create and Detail UI tests & Representative screenshots available \\
  Validation placement and state text & Client tests and validation captures & Representative evidence available \\
  Button hierarchy and disabled busy state & Client tests and UI contract & Dedicated submitting capture available \\
  Clipping, overlap, and horizontal overflow & E2E viewport assertions & Final-main-aligned E2E passed; human visual sign-off pending \\
  \bottomrule
\end{tabular}
\end{center}

The visual checklist is supported by the final-main-aligned E2E screenshot
inventory, the final-main rendered contract/repository/release/test evidence,
and the Project board capture. Human visual sign-off remains
explicitly open rather than being inferred from the release.

\end{document}
